\section{Related Work}

The starting point for this work is a tokenization inequality. Petrov et al.~\cite{petrov2023tokenization} showed that tokenizers trained on English-dominant corpora require up to 15$\times$ more tokens to represent the same content in non-English languages. For a Telugu speaker, this means every interaction with a commercial LLM is structurally more expensive---in cost, latency, and effective context window---before the model processes a single word. That asymmetry raises a deeper question: if the infrastructure already disadvantages non-English users, how well do LLMs actually serve them in practice?

For hundreds of millions of Indians, the practical answer involves romanized script. Rather than switching to a native-script keyboard, Telugu, Hindi, and Kannada speakers type their language in Latin characters---a behavior documented at 83\% prevalence for Hindi and 73\% for Telugu in real LLM deployments~\cite{khullar2025scriptgap}. This romanization is not standardized: the same word may be spelled a dozen ways depending on the speaker, a variability documented across South Asian languages in the Dakshina~\cite{roark2020dakshina} and Aksharantar~\cite{madhani2023aksharantar} corpora.

Prior work has begun to characterize what this means for LLM performance. Khullar et al.~\cite{khullar2025scriptgap} showed that romanized input degrades LLM performance on structured triage classification by up to 24 points relative to native script. Husain et al.~\cite{husain2024romansetu} took the complementary view---using romanization as an \emph{input strategy} to unlock multilingual capabilities---and found it improves performance on formal benchmarks. Both studies, however, evaluate against tasks with a known correct answer. Neither asks whether automatic metrics can detect quality failures on open-ended generation tasks, where no ground truth exists.

Standard evaluation benchmarks such as XQuAD~\cite{artetxe2020xquad}, IndicQA~\cite{doddapaneni2023indicqa}, and MLQA use formal native script and presuppose a reference answer to score against. Multilingual BERTScore~\cite{zhang2020bertscore} and BGE-M3 cosine similarity~\cite{chen2024bgem3} extend this paradigm to embedding space---but they still compare outputs to a reference, with no mechanism to detect whether a response uses the wrong register, chooses unnatural vocabulary, or reads like a translation rather than something a native speaker would write. We fill this gap: using native-speaker human evaluation as ground truth, we show that these metrics produce scores in a narrow, acceptable-looking band even when the outputs they are rating are rejected by the very speakers the system is supposed to serve.
